\section{Turing Machine --- State Renaming}
%%% generated by the blueprint pipeline from TCSlib.Complexity.TuringMachine.StateRenaming
%%%%%%%%%%%%%%%%%%%%%%%%%%%%%%%%%%%%%%%%%%%%%%%%%%%%%%%%%%%%%%%%%%%%%%%%%%%%%%%
\subsection*{Turing Machine --- State Renaming: Introduction}

This module transports the actions, configurations, and machines of the
multi-tape Turing machine model along maps of the state type: actions and
configurations can be renamed along an arbitrary function, while a whole
machine is relabelled along a bijection of state sets, and the relabelled
machine is shown to simulate the original step by step and, from the initial
configuration, at every time.

%%%%%%%%%%%%%%%%%%%%%%%%%%%%%%%%%%%%%%%%%%%%%%%%%%%%%%%%%%%%%%%%%%%%%%%%%%%%%%%
\subsection*{Turing Machine --- State Renaming: Results}

\begin{definition}[Renaming the states of an action]
\lean{Turing.Action.mapState}
\label{Turing.Action.mapState}
\leanok
\uses{Turing.Action}
Given a function $f$ from one state set to another, the one-step action of a
Turing machine with $k$ work tapes obtained from an action $a$ by applying
$f$ to its optional successor state.  The input-head movement, the
write-and-move instructions for the $k$ work tapes, and the optional output
symbol are unchanged, and a halting action (one with no successor state)
remains halting.
\end{definition}

\begin{definition}[Renaming the state of a configuration]
\lean{Turing.Cfg.mapState}
\label{Turing.Cfg.mapState}
\leanok
\uses{Turing.Cfg}
Given a function $f$ from one state set to another, the configuration of a
Turing machine with $k$ work tapes on a fixed input word obtained from a
configuration $c$ by applying $f$ to its optional state component.  The
input-head position, the work tapes and their head positions, and the output
are preserved, and a halted configuration (one with no state) stays halted.
\end{definition}

\begin{lemma}[State Renaming and Action Application]
\lean{Turing.Cfg.mapState_apply}
\label{Turing.Cfg.mapState_apply}
\leanok
\uses{Turing.Action, Turing.Action.apply, Turing.Action.mapState, Turing.Cfg, Turing.Cfg.mapState}
\difficulty{1}
Let $f$ be a function between state sets, let $a$ be a one-step action of a
$k$-work-tape Turing machine, and let $c$ be a configuration on a fixed input word.
Applying the action after renaming its successor state and the configuration state along
$f$ gives the same configuration as applying $a$ to $c$ and then renaming the resulting
state along $f$.
\end{lemma}

\begin{definition}[Transporting a machine along a state equivalence]
\lean{Turing.MultiTapeTM.relabelState}
\label{Turing.MultiTapeTM.relabelState}
\leanok
\uses{Turing.Action.mapState, Turing.MultiTapeTM}
Given a multi-tape Turing machine $M$ with $k$ work tapes and state set $S$,
and a bijection $e$ from $S$ to another state set $S'$, the machine over $S'$
obtained by relabelling the states of $M$ along $e$: its initial state is
$e(q_0)$, and its transition on a state $q'$, the scanned input symbol, and
the $k$ scanned work-tape symbols is the action that $M$ prescribes in the
state $e^{-1}(q')$ for the same scanned symbols, with its successor state
renamed along $e$.  A bijection, rather than an arbitrary function, is
required here: a map identifying two states with different transitions would
leave no well-defined transition table.
\end{definition}

\begin{lemma}[State relabelling commutes with one step]
\lean{Turing.MultiTapeTM.relabelState_step}
\label{Turing.MultiTapeTM.relabelState_step}
\leanok
\uses{Turing.Cfg, Turing.Cfg.inputSymbol, Turing.Cfg.mapState, Turing.Cfg.mapState_apply, Turing.Cfg.workTapeSymbols, Turing.MultiTapeTM, Turing.MultiTapeTM.relabelState, Turing.MultiTapeTM.step}
\difficulty{4}
Let $M$ be a multi-tape Turing machine with $k$ work tapes, let $e$ be a bijection from
its state set to another state set, let $x$ be an input word, and let $c$ be a
configuration of $M$ on $x$. Taking one step of the relabelled machine from the
configuration obtained by applying $e$ to the state of $c$ gives the same configuration
as taking one step of $M$ from $c$ and then applying $e$ to its state.
\end{lemma}

\begin{lemma}[State relabelling commutes with a run from the initial configuration]
\lean{Turing.MultiTapeTM.relabelState_runFrom_init}
\label{Turing.MultiTapeTM.relabelState_runFrom_init}
\leanok
\uses{Turing.Cfg.mapState, Turing.MultiTapeTM, Turing.MultiTapeTM.initCfg, Turing.MultiTapeTM.relabelState, Turing.MultiTapeTM.relabelState_step, Turing.MultiTapeTM.runFrom, Turing.MultiTapeTM.runFrom_comm_of_step}
\difficulty{1}
Let $tm$ be a multi-tape Turing machine with state set $S$, let $e:S\simeq S'$ be a
bijection, let $x$ be an input word, and let $t\in\mathbb{N}$. After $t$ steps from
its initial configuration on $x$, the machine obtained by relabelling the states of $tm$
along $e$ has the same configuration as the configuration reached by $tm$ after $t$
steps from its own initial configuration, with the state relabelled by $e$.
\end{lemma}

%%% added by the blueprint pipeline
\subsection*{Turing Machine --- State Renaming: Further results}

\begin{theorem}[Runs agree under an embedding on guarded states]
\lean{Turing.MultiTapeTM.runFrom_mapState_of_agreeOn}
\label{Turing.MultiTapeTM.runFrom_mapState_of_agreeOn}
\leanok
\uses{Turing.Action.apply, Turing.Action.mapState, Turing.Cfg, Turing.Cfg.inputSymbol, Turing.Cfg.mapState, Turing.Cfg.mapState_apply, Turing.Cfg.workTapeSymbols, Turing.MultiTapeTM, Turing.MultiTapeTM.AgreeOn, Turing.MultiTapeTM.runFrom, Turing.MultiTapeTM.runFrom_comm_of_step, Turing.MultiTapeTM.runFrom_eq_of_agreeOn, Turing.MultiTapeTM.step}
\difficulty{5}
Let the source and host be multi-tape Turing machines with $k$ work tapes, a common
symbol alphabet, and state sets $S$ and $H$, and let $x$ be an input word. Let
$e:S\hookrightarrow H$ be an injective map and let $G$ be a predicate on source states.
Suppose that, for every $q\in S$ satisfying $G(q)$ and every scanned input symbol and
tuple of work-tape symbols, the host transition at $e(q)$ equals the source transition
at $q$ with its successor state mapped by $e$. For any source configuration $c$ on $x$
and any $t\in\mathbb{N}$, if at every time $u<t$ the run of the source machine from $c$
has state $q$ only when $G(q)$ holds, then running the host for $t$ steps from the
configuration obtained by mapping the state of $c$ along $e$ yields the configuration
reached by the source after $t$ steps, with its state mapped along $e$.
\end{theorem}

\begin{theorem}[State relabelling preserves visited cells]
\lean{Turing.MultiTapeTM.relabelState_visitedByTapeHead}
\label{Turing.MultiTapeTM.relabelState_visitedByTapeHead}
\leanok
\uses{Turing.Cfg.mapState, Turing.MultiTapeTM, Turing.MultiTapeTM.initCfg, Turing.MultiTapeTM.relabelState, Turing.MultiTapeTM.relabelState_runFrom_init, Turing.MultiTapeTM.visitedByTapeHead}
\difficulty{1}
Let $M$ be a multi-tape Turing machine with $k$ work tapes, let $e$ be a bijection from
its state set to another state set, let $x$ be an input word, let $t\in\mathbb{N}$, and
choose a work tape $i<k$. The set of integer positions visited by the head of tape $i$
from the initial configuration on $x$ through step $t$ is the same for $M$ and the
machine obtained by relabelling its states along $e$.
\end{theorem}

\begin{theorem}[State relabelling preserves space use]
\lean{Turing.MultiTapeTM.relabelState_spaceUsed}
\label{Turing.MultiTapeTM.relabelState_spaceUsed}
\leanok
\uses{Turing.MultiTapeTM, Turing.MultiTapeTM.initCfg, Turing.MultiTapeTM.relabelState, Turing.MultiTapeTM.relabelState_visitedByTapeHead, Turing.MultiTapeTM.spaceUsed, Turing.MultiTapeTM.spaceUsedByTape}
\difficulty{1}
Let $M$ be a multi-tape Turing machine with $k$ work tapes, let $e$ be a bijection from
its state set to another state set, let $x$ be an input word, and let $t\in\mathbb{N}$.
The total number of distinct work-tape cells visited from the initial configuration on
$x$ through step $t$ is the same for $M$ and the machine obtained by relabelling its
states along $e$.
\end{theorem}
